%======================================================================
% CAPÍTULO 11 - MARCO TEÓRICO
%======================================================================
\chapter{Marco Teórico}

\dropcap{E}{l} desarrollo de sistemas de nómina de alto rendimiento se fundamenta en principios de la ingeniería de sistemas, la computación de alto rendimiento y las matemáticas financieras. Este capítulo establece las bases teóricas que sustentan el diseño de Sentinel.

\section{Computación en Memoria}

La computación en memoria representa un paradigma donde los datos se procesan completamente en la memoria RAM del sistema, eliminando la latencia asociada a las operaciones de entrada/salida contra bases de datos tradicionales.

\begin{center}
\begin{tabular}{lcc}
\toprule
\textbf{Tipo de Acceso} & \textbf{Latencia} & \textbf{Orden} \\
\midrule
Memoria RAM & 100 ns & $10^{-7}$ s \\
Disco SSD & 100 $\mu$s & $10^{-4}$ s \\
Disco Mecánico & 10 ms & $10^{-2}$ s \\
\bottomrule
\end{tabular}
\end{center}

Esta diferencia de hasta cinco órdenes de magnitud justifica el enfoque de Sentinel:

\begin{itemize}
    \item \textbf{Sin consultas SQL}: Durante el procesamiento no hay acceso a base de datos
    \item \textbf{Búsquedas O(1)}: Mediante tablas hash en memoria
    \item \textbf{Procesamiento paralelo}: Todos los núcleos de CPU trabajan simultáneamente
    \item \textbf{Tiempos reducidos}: De horas a segundos en el procesamiento
\end{itemize}

\section{Arquitectura Hexagonal}

La arquitectura hexagonal, también conocida como \textit{Ports and Adapters}, es un patrón de diseño que busca desacoplar la lógica de negocio de las dependencias externas.

\textbf{Ports:} Interfaces abstractas que definen cómo el núcleo de la aplicación interactúa con el mundo exterior.

\textbf{Adapters:} Implementaciones concretas de los ports que conectan con tecnologías específicas (gRPC, archivos, bases de datos).

En Sentinel el desacoplamiento se manifiesta así:

\begin{itemize}
    \item \textbf{Núcleo}: Motor de cálculo, lógica de fusión, algoritmos de distribución
    \item \textbf{Ports de entrada}: Manifiestos JSON, streams gRPC
    \item \textbf{Ports de salida}: Archivos CSV, logs de auditoría
    \item \textbf{Adapters}: Cliente gRPC, exportador CSV
\end{itemize}

\section{Paralelismo de Datos con Rayon}

Rayon es una biblioteca de Rust que proporciona paralelismo de datos con una API mínima. Utiliza un modelo de \textit{work stealing} que distribuye automáticamente el trabajo entre los núcleos disponibles.

Para un conjunto de datos de tamaño $N$ procesado en $P$ núcleos, el speedup teórico máximo es $P$. En la práctica, el speedup está limitado por la fracción secuencial según la \textbf{Ley de Amdahl}:

\[
S = \frac{1}{(1 - f) + \frac{f}{P}}
\]

donde $f$ es la fracción paralelizable del algoritmo.

Sentinel utiliza Rayon para paralelizar:

\begin{itemize}
    \item Cálculo de nómina por beneficiario
    \item Evaluación de fórmulas de primas
    \item Generación de archivos de exportación
\end{itemize}

\section{Precisión en Cálculos Monetarios}

Los cálculos financieros requieren precisión absoluta. El estándar IEEE 754 para números de punto flotante puede introducir errores de redondeo que, acumulados en millones de registros, generan diferencias significativas.

\textbf{Problema:} $0.1 + 0.2 \neq 0.3$ en punto flotante (produce $0.30000000000000004$)

\textbf{Solución:} Sentinel usa aritmética de enteros. Todos los valores monetarios se convierten a centavos:

\[
\text{valor\_centavos} = \lfloor \text{valor\_decimal} \times 100 + 0.5 \rfloor
\]

Las operaciones se realizan con enteros de 64 bits (\texttt{i64}) y el resultado final se reconvierte dividiendo por 100.

\section{Pipeline Asíncrono}

El procesamiento asíncrono permite solapamiento de operaciones de E/S con procesamiento de CPU, maximizando la utilización de recursos.

En un pipeline optimizado:
\begin{itemize}
    \item Mientras la CPU procesa el lote $N$
    \item La red descarga el lote $N+1$
    \item El disco escribe el lote $N-1$
\end{itemize}

Sentinel implementa este patrón mediante Tokio, utilizando streams asíncronos para consumo de datos gRPC, spawn de tareas para procesamiento paralelo, y \textit{back-pressure} para evitar sobrecarga de memoria.

\vfill
\newpage
